% 解説プリントのコメント（saiten が書く。./tools/kaisetsu.ps1 がこのファイルからプリントを作る）
% 2026-11-10 スキャン（2026-11-10_05_46_36_PM_からスキャン_p1/p2）の採点
%% matome
\begin{daiji}
  \item 「示せ」「確かめよ」と問われた式は、結果を使う前に1行でも導く。$\hat{U}^{\dagger}=\ee^{\ii\hat{H}t/\hbar}$ は $\hat{H}^{\dagger}=\hat{H}$ から、初期条件は $\hat{U}(0)=\hat{1}$ から出る（問題1 (1)(3)）。
  \item 期待値は数であり、最後の行がケットになっていないかを確かめる。$\expval{\hat{H}}_{t}=(E_{1}+2E_{2})/3$ であって $\frac{1}{3}\ket{E_{1}}+\frac{2}{3}\ket{E_{2}}$ ではない（問題1 (2)）。
  \item 数値の答えには単位を付ける（問題3 (2)）。
\end{daiji}
\begin{maemo}
  \item 単位の誤り・書き落とし：第6章 問題6 (5)（2026-11-01）で $\delta p$ の単位を $\unit{\metre}$ と書いた。今回は問題3 (2) で単位が無い。数値を出したら、その行のうちに単位を書き、次元が問いの量と合うかを見る。
  \item 条件が成り立つことの確認を書かない：第7章 問題3 (1)（2026-11-04）で $\varphi'(L)=\varphi'(0)$ が自動で成り立つことを書かなかった。今回は問題1 (1) の初期条件。「満たす」と書く前に、$t=0$ などを代入した式を1行置く。
\end{maemo}
%% mondai 1
\begin{kotaeran}
  \machigai{1}{○}{記述} 級数を項別に微分して $\ii\hbar\,\dd\hat{U}/\dd t=\hat{H}\hat{U}$ を出すところは正しい。最後に「シュレーディンガー方程式と初期条件を満たす」と書いたが、初期条件の根拠が式にも文にも無い。$t=0$ では $n\ge1$ の項が消えて $\hat{U}(0)=\hat{1}$、よって $\ket{\psi(t)}$ は $t=0$ で $\ket{\psi(0)}$ に戻る、を書く。
  \machigai{2}{○}{記述} $\ket{\psi(t)}$ と $\braket{\psi(t)}{\psi(t)}=1$ は正しく、$\expval{\hat{H}}_{t}$ も途中の $\frac{1}{3}E_{1}\braket{E_{1}}{E_{1}}+\frac{2}{3}E_{2}\braket{E_{2}}{E_{2}}$ までは正しい。最後の行を $\frac{1}{3}\ket{E_{1}}+\frac{2}{3}\ket{E_{2}}$ とケットで書いている。$\braket{E_{j}}{E_{j}}=1$ を入れると、正しくは $\expval{\hat{H}}_{t}=(E_{1}+2E_{2})/3$（数）。
  \machigai{3}{○}{記述} $\hat{U}^{\dagger}\hat{U}=\exp(-\ii\hat{H}t/\hbar+\ii\hat{H}t/\hbar)=\hat{1}$ とノルムの保存は正しい。問われている $\hat{U}(t)^{\dagger}=\exp(\ii\hat{H}t/\hbar)$ を示さずに使っている。級数の各項で $(c\hat{H}^{n})^{\dagger}=c^{*}\hat{H}^{n}$（$\hat{H}^{\dagger}=\hat{H}$ による）を使って導く。
\end{kotaeran}
\begin{rikai}
  \item $\hat{U}^{\dagger}=\ee^{\ii\hat{H}t/\hbar}$ は $\hat{H}$ がエルミートだから成り立つ。$\hat{H}$ がエルミートでなければ $\hat{U}^{\dagger}\hat{U}=\hat{1}$ は言えず、ノルムも保存しない。つまり (3) の要点はこの1行にある。
  \item 時間発展の解は「方程式」と「初期条件」の2つを満たして初めて決まる。$\hat{U}(0)=\hat{1}$ は級数に $t=0$ を入れるだけで出る。
  \item $\matrixel{\psi}{\hat{A}}{\psi}$ はブラとケットで挟んだ数である。最後の行に $\ket{\ }$ が残っていたら誤り。
\end{rikai}
%% mondai 3
\begin{kotaeran}
  \machigai{2}{○}{単位} $\delta\expval{x}=-0.196$ と単位が無い。数値は正しく、質量にも $\hbar$ にもよらないことも (1) の式から読み取れる。正しくは $\qty{-0.196}{\metre}$（$\qty{0.196}{\metre}$ 下がる）。
\end{kotaeran}
\begin{rikai}
  \item 数値の答えには単位を付ける。$gt^{2}$ は $\unit{\metre\per\second\squared}\times\unit{\second\squared}=\unit{\metre}$ で、位置の変化の次元になっていることを確かめる。
\end{rikai}
